\color{unimain}{\subsection{Внешнее отключение трансформатора}}
\color{black}

\begin{enumerate}

\item
Внешнее отключение трансформатора предназначено для приема внешних сигналов и формирования требуемого действия на отключение выключателей трансформатора.

\item
Алгоритм работы логики внешнего отключения трансформатора приведен на рисунке \ref{pict:VO}, уставки функции приведены в таблице \ref{app_set:vo}.

\begin{figure}[H]
\centering
\includegraphics[width=0.75\textwidth,height=0.75\textheight,keepaspectratio]{\fbpath/07.01 РЗТ/011.1761 ВО/img/VO.pdf}
\captionof{figure}{Функциональная схема ВО}\label{pict:VO}
\end{figure}

\item
Предусматривается возможность приема внешних сигналов \textit{«Откл. от ДЗО ВН»}, \textit{«Откл. от ДЗО СН»}, \textit{«Откл. от ДЗО НН»}, \textit{«Откл. от УРОВ В1 НН1(НН)»}, \textit{«Откл. от УРОВ В2 НН1(НН)»}, \textit{«Откл. от УРОВ В1 НН2(СН)»} и \textit{«Откл. от УРОВ В2 НН2(СН)»} на отключение трансформатора с пуском УРОВ и запретом АПВ, а также прием трех дополнительных внешних сигналов на отключение \textit{«Внеш. откл1»}, \textit{«Внеш. откл2»} и \textit{«Внеш. откл3»} с возможностью контроля по току. При необходимости ввести контроль по току для сигналов \textit{«Откл. от ДЗО ВН»}, \textit{«Откл. от ДЗО СН»}, \textit{«Откл. от ДЗО НН»}, \textit{«Откл. от УРОВ В1 НН1(НН)»}, \textit{«Откл. от УРОВ В2 НН1(НН)»}, \textit{«Откл. от УРОВ В1 НН2(СН)»} и \textit{«Откл. от УРОВ В2 НН2(СН)»} их можно сконфигурировать на дополнительные сигналы внешнего отключения.

\item
Порог срабатывания контроля по току ВО задается уставками \textbf{«Iср ВО-1 ВН»}, \textbf{«Iср ВО-1 СН»}, \textbf{«Iср ВО-2 ВН»}, \textbf{«Iср ВО-2 СН»}, \textbf{«Iср ВО-3 ВН»} и \textbf{«Iср ВО-3 СН»}, а выдержки времени срабатывания ВО регулируются уставками \textbf{«Тср ВО-1»}, \textbf{«Тср ВО-2»} и \textbf{«Тср ВО-3»}.  Ввод в работу ВО осуществляется с помощью программных переключателей \textbf{«ВО-1»}, \textbf{«ВО-2»} и \textbf{«ВО-3»}.

\item
Оперативный ввод/вывод ВО производится с помощью БУ (блоков управления функциями) \textit{«Вывод ВО-1»}, \textit{«Вывод ВО-2»} и \textit{«Вывод ВО-3»}.

\item
Возможность контроля ВО по токам ВН и СН задается программными переключателями \textbf{«Контроль ВО-1»}, \textbf{«Контроль ВО-2»} и \textbf{«Контроль ВО-3»}, имеющими четыре положения:
\begin{itemize}
\item \textit{«Выведено»} -- ВО без контроля по току;
\item \textit{«Iф ВН»} -- ВО с контролем максимально фазного тока ВН;
\item \textit{«Iф СН»} -- ВО с контролем максимально фазного тока СН;
\item \textit{«Iф ВН, СН»} -- ВО с контролем максимально фазного тока ВН или СН.
\end{itemize}

\item
Режим действия внешнего отключения на выключатели трансформатора задается с помощью программных переключателей \textbf{«Режим ВО-1»}, \textbf{«Режим ВО-2»} и \textbf{«Режим ВО-3»}, имеющих два положения:
\begin{itemize}
\item \textit{«Без АПВ»} – ВО с запретом АПВ;
\item \textit{«С АПВ»} – ВО без запрета АПВ.
\end{itemize}

\item
Предусмотрен контроль исправности внешних дополнительных сигналов на отключения. При наличии сигнала внешнего отключения и отсутствии разращения по току (при этом контроль внешнего отключения по токам ВН или СН должен быть введен) более 5 секунд, формируется сигнал о неисправности соответствующего внешнего отключения.

\end{enumerate} 